\documentclass[12pt]{article}
\usepackage{amsmath,amssymb,amsfonts,amsthm}
\usepackage[margin=1in]{geometry}

\begin{document}

\begin{center}
{\Large\bfseries Chapter 9, Problem 1}\\[6pt]
Byron \& Fuller, \textit{Mathematics of Classical and Quantum Physics}
\end{center}

\bigskip

\textbf{Problem.} If $A$ is a completely continuous Hermitian operator with only a finite number of nonzero eigenvalues, show that the eigenvectors of $A$ can be chosen to form a complete orthonormal set.

\bigskip
\textbf{Solution.}

Since $A$ is a completely continuous (compact) Hermitian operator on a Hilbert space $H$, by Theorem 9.4, $A$ has at most a countable number of eigenvalues, and any nonzero eigenvalue has finite multiplicity. We are told that there are only a \emph{finite} number of nonzero eigenvalues, say $\lambda_1, \lambda_2, \ldots, \lambda_N$ (counted with multiplicity).

For each nonzero eigenvalue $\lambda_k$, let $H_k$ denote the corresponding eigenspace. Since each eigenspace is finite-dimensional, we can choose an orthonormal basis for each $H_k$ by the Gram--Schmidt process. Eigenvectors corresponding to \emph{different} eigenvalues of a Hermitian operator are automatically orthogonal, so the union of all these orthonormal bases,
\[
\{x_1, x_2, \ldots, x_M\},
\]
forms an orthonormal set spanning the direct sum $H_1 \oplus H_2 \oplus \cdots \oplus H_N$.

Now consider the orthogonal complement of this span:
\[
H_0 = \bigl(\text{span}\{x_1, \ldots, x_M\}\bigr)^\perp.
\]
For any $y \in H_0$, we have $(y, x_k) = 0$ for all $k$. We claim $Ay = 0$ for all $y \in H_0$.

Since $A$ is Hermitian, $H_0$ is an invariant subspace of $A$ (if $y \perp x_k$, then $(Ay, x_k) = (y, Ax_k) = \lambda_k(y, x_k) = 0$, so $Ay \in H_0$). The restriction $A|_{H_0}$ is also a completely continuous Hermitian operator. If $A|_{H_0}$ had any nonzero eigenvalue, it would be a nonzero eigenvalue of $A$ itself, contradicting the fact that all nonzero eigenvalues of $A$ have their eigenvectors in $H_1 \oplus \cdots \oplus H_N$. Therefore $A|_{H_0}$ has no nonzero eigenvalues, and by Theorem 9.2 (the norm of a completely continuous self-adjoint operator equals its largest eigenvalue in absolute value), we conclude $\|A|_{H_0}\| = 0$, hence $Ay = 0$ for all $y \in H_0$.

Thus every vector in $H_0$ is an eigenvector of $A$ with eigenvalue $0$. We can choose an orthonormal basis $\{e_\alpha\}$ for $H_0$ (using Zorn's lemma or, if $H_0$ is separable, the Gram--Schmidt process on a countable dense subset).

The combined set $\{x_1, \ldots, x_M\} \cup \{e_\alpha\}$ is an orthonormal set of eigenvectors of $A$ (with eigenvalues $\lambda_k$ for the $x_k$ and $0$ for the $e_\alpha$). Since $H = (H_1 \oplus \cdots \oplus H_N) \oplus H_0$, this set is a \emph{complete} orthonormal set (basis) for $H$.

\hfill $\blacksquare$

\end{document}
